%% kapa: exam=1 round=29 year=2018 period=1 form=A subject=civil_law subject_ko=민법 seq=6 number=6
%% source: Q-Net 감정평가사 기출문제 (cst00309) artlSeq=5202808 file=extracted/1008/5202808/2018년 제29회 감정평가사 1차시험 기출문제/2018년 제29회 감정평가사 1차시험 기출문제/1교시 A형.pdf
%% conversion: kordoc 4.24.51 -> tools/gen1.py (auto); review=proofread-sonnet; flags=none
\begin{question}{6}
민법상 법인의 정관에 관한 설명으로 옳은 것은?
\begin{choices}
  \item 사단법인의 정관변경은 법원의 허가를 얻지 않으면 그 효력이 없다.
  \item 사단법인에서 이사의 대표권에 대한 제한은 정관에 기재되지 않더라도 효력이 있다.
  \item 재단법인 설립자는 정관에 그 존립시기나 해산사유를 기재하고 기명날인하여야 한다.
  \item 재단법인 설립자가 이사의 임면방법을 정하지 아니하고 사망한 경우, 이해관계인의 청구에 의하여 주무관청이 이를 정한다.
  \item 재단법인의 재산보전을 위하여 적당한 때에는 정관에 변경방법이 없더라도 명칭 또는 사무소의 소재지를 변경할 수 있다.
\end{choices}
\end{question}
\answer{5}
